This is a small CPU study of one analytic system. (i) Single element, pair-only potential, at most 13 atoms; the near-linear behaviour of symmetry-function KRR and the descriptor ranking may not transfer to many-body or multi-element systems. (ii) The minima library (one to four minima per size) is shared by train and test, so we measure interpolation around known minima, not generalisation to new structures or sizes. (iii) Energies are truncated at $E\le50$, which excludes the strongly repulsive region. (iv) 3--5 seeds; tests with 3 seeds have little power, are uncorrected for multiplicity, and the H3 verdict rests on $p$-values of 0.012 and 0.032. (v) Hyperparameters: the KRR grid was widened after an audit of the first version had shown the test-set consequences of the narrow grid, so the grid is not a pre-registered choice, although its bounds were set on validation error; the selected $\lambda$ of symmetry-function KRR sits on the smallest numerically feasible value in 13 of 14 runs, and such nearly singular solves may be sensitive to the linear-algebra library. (vi) All baselines are our reimplementations; the MLP architecture and step count were not tuned per system, the step-count check shows the sorted-distance MLP was not converged, and the poor symmetry-function networks say more about our training setup than about Behler--Parrinello networks. (vii) Raw-coordinate models were trained on randomly rotated, randomly ordered data, so they never had a chance to fit; training on canonically aligned data was not tried, and the H2 refutation applies to this setting only. (viii) Test sets differ between training sizes, so learning-curve slopes include test-sampling noise. (ix) Energies only: no forces, no dynamics, no per-$N$ breakdown. (x) Invariance is measured with one random transform per test configuration; the exact zeros of two MLP systems may reflect their single-precision inputs. A fuller study would use unseen minima and sizes, forces, more seeds, tuned networks, SOAP-type kernels and equivariant networks.
